



\label{sec:fidelity estimation}









We provide a quantum algorithm for estimating the square root fidelity $\sqrt{\rF}(\rho^\sA, \sigma^\sA)$ and evaluate the query and sample complexities in the three models.
In the sample access models, the sample cost depends on whether collective operations over multiple samples are allowed or not. We also provide a lower bound on the query and sample complexities for fidelity estimation, which in turn implies a corresponding lower bound for the Uhlmann transformation. The results are organized as follows:
\begin{itemize}
    \item The purified query access model in \Cref{sec:fidelity estimation PQA}.
    \item The purified and mixed sample access models in \Cref{sec:fidelity estimation purif and mixed sample}: collective case in \Cref{sec:fidelity estimation collective} and non-collective case in Secs.~\ref{sec:fidelity estimation PSS} and~\ref{sec:fidelity estimation SS}.
    \item A lower bound on the query and sample complexities in \Cref{sec:lower bound}.
\end{itemize}





The key idea of our approach is first to apply the Uhlmann transformation $V^\sB$ to the state $\ket{\rho}^{\sA\sB}$, and then, to estimate the square root fidelity between two pure states $V^\sB\ket{\rho}^{\sA\sB}$ and $\ket{\sigma}^{\sA\sB}$.
Due to the Uhlmann's theorem (Theorem~\ref{thm:Uhlmann theorem}), $\sqrt{\rF}(V^\sB\ket{\rho}^{\sA\sB}, \ket{\sigma}^{\sA\sB}) = \big|\bra{\sigma}V^\sB\ket{\rho}^{\sA\sB}\big|$ coincides with $\sqrt{\rF}(\rho^\sA, \sigma^\sA)$.
It is important to note that, when we implement the Uhlmann transformation using the quantum algorithms, $V^\sB$ can be realized approximately.
We need to evaluate how the approximation error of $V^\sB$ affects the fidelity estimation.


Comparing our approach with previous ones, a notable distinction is that previous approaches mainly aimed to directly prepare the square root of the states and discarded the ``purifying system'' midway, even when purified access was available.
In contrast, our approach keeps the purifying system throughout the process, enabling us to employ a square root fidelity estimation algorithm for pure states.
While fidelity estimation between mixed states is generally difficult~\cite{gilyen2022improvedfidelity, Rethinasamy2023estimatingdist}, fidelity estimation between pure states is more feasible, and optimal algorithms are known~\cite{Wang2024OptimalFidelityPure, wang2024sampleoptimal,Fang2025OptimalFidelityToPure}.
This results in a substantial speedup for our algorithm in the purified access model.












%=====================================================================================================================================================================================================



\subsection{In the purified query access model}
\label{sec:fidelity estimation PQA}


Our goal is to estimate the fidelity $\sqrt{\rF}(\rho^\sA, \sigma^\sA)$ using the Uhlmann transformation algorithm in the purified access model.
For simplicity, we sometimes abbreviate $\sqrt{\rF}(\rho^\sA, \sigma^\sA)$ as $\sqrt{\rF}$.
We propose a quantum query algorithm $\mathtt{UhlFidelityEstPurifQuery}$, whose performance is 
characterized by the following theorem.

\begin{theorem}[Square root fidelity estimation in the purified query access model]
\label{thm:fid est quel}
Let $\delta \in (0, 1)$. Then, the quantum query algorithm $\mathtt{UhlFidelityEstPurifQuery}(U_\rho^{\sA\sB}, U_\sigma^{\sA\sB}; \delta)$,
which is given by \Cref{alg:fidelestquery}, outputs $\sqrt{\til{\rF}}$ such that $\big|\sqrt{\til{\rF}} - \sqrt{\rF}\big| \leq \delta$ with probability at least $2/3$, using $q = \cO\big(\f{1}{\delta}{\rm min}\big\{\f{1}{s_{\rm min}}, \f{r}{\delta}\big\}\log{\big(\f{1}{\delta}\big)}\big)$ queries to $U_\rho^{\sA\sB}$, $U_\sigma^{\sA\sB}$, and their inverses. 
The quantum circuit of this algorithm consists of $\cO\big(q\log{(d_\sA d_\sB)} + (\log{(1/\delta}))^2\big)$ one- and two-qubit gates, and $\cO(\log{(d_\sA d_\sB)} + \log{(1/\delta}))$ qubits suffice at any one time.
\end{theorem}


\begin{algorithm}[h]
\caption{Square root fidelity estimation algorithm in the purified query access model \\ \parbox{\linewidth}{\centering $\mathtt{UhlFidelityEstPurifQuery}(U_\rho^{\sA\sB}, U_\sigma^{\sA\sB}; \delta)$ (In Theorem~\ref{thm:fid est quel})}}
\label{alg:fidelestquery}
\SetKwInput{KwInput}{Input}
\SetKwInput{KwOutput}{Output}
\SetKwInput{KwParameters}{Parameters}

\SetAlgoNoEnd
\SetAlgoNoLine
\KwInput{Unitary oracles $U_\rho^{\sA\sB}$, $U_\sigma^{\sA\sB}$, and their inverses.}
\KwParameters{$\delta \in (0, 1)$.}
\KwOutput{Real number $\sqrt{\til{\rF}}$.}
\SetAlgoLined


Set $\til{W}_\tF^{\sB\sD} \gets \mathtt{UhlmannPurifiedQuery}(U_\rho^{\sA\sB}$, $U_\sigma^{\sA\sB}; \delta, \tF)$ (Algorithm~\ref{alg:Uhl purif query}). \\
Set $\sqrt{\til{\rF}} \gets \mathtt{SqrtAmpEstQuery}(\til{W}_\tF^{\sB\sD}U_\rho^{\sA\sB}, U_\sigma^{\sA\sB}; \delta/2)$ (Lemma~\ref{lem:square root amp est query}). \\
Return $\sqrt{\til{\rF}}$.
\end{algorithm}


In \Cref{alg:fidelestquery}, we use the \emph{square root amplitude estimation} algorithm $\mathtt{SqrtAmpEstQuery}$~\cite{Wang2024OptimalFidelityPure}, which is known to be optimal in terms of the number of queries in the case where both states are pure.
Here, we use a slightly adjusted version of the result from Ref.~\cite{Wang2024OptimalFidelityPure}.

\begin{theorem}[Square root amplitude estimation in the purified query access model~{\cite[Theorem III.4]{Wang2024OptimalFidelityPure}}]
\label{lem:square root amp est query}
Suppose that $W$ and $U$ are unitaries such that $W\ket{0} = \sqrt{c}U\ket{0} + \sqrt{1- c}\ket{\hspace{-1mm}\perp}$,
where $\ket{\hspace{-2mm}\perp}$ is a state satisfying $\bra{\perp\hspace{-2mm}}U\ket{0} = 0$.
Then, for any $\delta \in (0, 1)$, there is a quantum query algorithm $\mathtt{SqrtAmpEstQuery}(W, U; \delta)$ that outputs $\sqrt{\til{c}}$ such that $\big|\sqrt{\til{c}} - \sqrt{c}\big| \leq \delta$ holds with probability at least $2/3$, with $\cO(1/\delta)$ queries to $W$, $U$, and their inverses.
The quantum circuit of this algorithm consists of $\cO\big((\log{(1/\delta)})^2\big)$ one- and two-qubit gates, and $\cO(\log{(1/\delta)})$ qubits suffice at any one time.
\end{theorem}

The algorithm $\mathtt{SqrtAmpEstQuery}$ is based on the \emph{quantum phase estimation}.
To estimate a phase of a target unitary $Q$ with accuracy $\delta$ and success probability $p_{\rm succ} \geq 1 - \eta$, the quantum phase estimation uses a total of $g = \cO(1/(\eta\delta))$ applications of $Q$. The quantum circuit includes $\cO\big((\log(1/\delta))^2\big)$ one- and two-qubit gates, which come from the quantum Fourier transform part, and uses $\cO(\log(1/\delta))$ auxiliary qubits.
Here, we took the overlap between the input state to the quantum phase estimation and the eigenstate of $Q$ corresponding to the estimated phase to be of constant order. In fact, while the success probability depends on the inverse of the overlap, the probability can be amplified with only a constant number of repetitions of $Q$ when the overlap is of constant order. 


In the fidelity estimation, the target unitary is given by a product of reflection unitaries: 
\begin{align}
    Q&= e^{i\pi W\ketbra{0}{0}W^\dag}e^{i\pi U\ketbra{0}{0}U^\dag} \\
    &= (\bI - 2W\ketbra{0}{0}W^\dag)(\bI - 2U\ketbra{0}{0}U^\dag)  \\
    \label{inteq:151}
    &= We^{i\pi\ketbra{0}{0}}W^\dag U e^{i\pi\ketbra{0}{0}}U^\dag.
\end{align}
The input state is taken as $W\ket{0}$, since this state has a constant overlap with eigenstates of $Q$.
Then, it was shown that performing quantum phase estimation on $Q$ and a simple post-processing calculation yields an outcome that sufficiently approximates the square root fidelity between two pure states with high probability~\cite{Wang2024OptimalFidelityPure}.

To perform the quantum phase estimation for $Q$ in Eq.~\eqref{inteq:151}, we need the control version of $e^{i\pi\ketbra{0}{0}} = \bI - 2\ketbra{0}{0}$, in addition to $W$ and $U$. 
It can be easily constructed using multi-controlled NOT gates that are implemented with a polynomial number of one- and two-qubit gates in the number of qubits~\cite{nielsen2010quantum, saeedi2013NoAncilla, yoder2014fixed}.






With the above discussion, we now prove Theorem~\ref{thm:fid est quel}.

\begin{proof}[Proof of Theorem~\ref{thm:fid est quel}]





As discussed in \Cref{sec:uhlfidquery} (in particular, from just above Eq.~\eqref{inteq:14} to Eq.~\eqref{inteq:13}), a unitary $\til{W}_\tF^{\sB\sD} = \mathtt{UhlmannPurifiedQuery}(U_\rho^{\sA\sB}, U_\sigma^{\sA\sB}; \delta, \tF)$, which is an exact block-encoding of $P_\sgn(M^\sB)$, satisfies that 
\begin{align}
\label{inteq:53}
    \big|\sqrt{\rF}\big(\til{W}_\tF^{\sB\sD}\ket{\rho}^{\sA\sB}\ket{0}^\sD, \ket{\sigma}^{\sA\sB}\ket{0}^\sD\big) - \sqrt{\rF}\big| \leq \delta/2, 
\end{align}
where $M^\sB = \tr_{\hat{\sA}}\big[\ketbra{\sigma}{\rho}^{\hat{\sA}\sB}\big]$.
The number of queries for implementing $\til{W}_\tF^{\sB\sD}$ is given by $\cO\big(\log{(1/\delta)}/\beta_\tF\big)$, where $\beta_\tF = \max\{s_{\rm min}, \delta/(8r)\}$.
Remember that $s_{\rm min}$ and $r$ are the minimum non-zero singular value and the rank of $\sqrt{\sigma^\sA}\sqrt{\rho^\sA}$, respectively.



Next, we consider estimating the value of $\sqrt{\rF'} = \sqrt{\rF}\big(\til{W}_\tF^{\sB\sD}\ket{\rho}^{\sA\sB}\ket{0}^\sD, \ket{\sigma}^{\sA\sB}\ket{0}^\sD\big)$. 
Using the Uhlmann transformation algorithm, we can obtain the states $\til{W}_\tF^{\sB\sD}\ket{\rho}^{\sA\sB}\ket{0}^{\sD}$ such that
\begin{align}
    \label{inteq:170}
    \til{W}_\tF^{\sB\sD}\ket{\rho}^{\sA\sB}\ket{0}^{\sD}
    &= \til{W}_\tF^{\sB\sD}U_\rho^{\sA\sB}\ket{0}^{\sA\sB\sD}    \\
    &= \sqrt{\rF'}e^{i\theta}\ket{\sigma}^{\sA\sB}\ket{0}^{\sD} + \ket{\perp~}^{\sA\sB\sD}, \\
    &= \sqrt{\rF'}U_\sigma^{\sA\sB}\ket{0}^{\sA\sB\sD} + \ket{\perp~}^{\sA\sB\sD},
\end{align}
where $\ket{\perp~}^{\sA\sB\sD}$ is a state satisfying $\bra{\perp\!}U_\sigma^{\sA\sB}\ket{0}^{\sA\sB\sD} = 0$.
A certain phase $\theta$ does not affect the result and is hence absorbed into $U_\sigma^{\sA\sB}$.
Thus, from \Cref{lem:square root amp est query}, when we perform the quantum phase estimation for a target unitary
\begin{align}
\label{inteq:48}
    Q^{\sA\sB\sD} = \til{W}_\tF^{\sB\sD}U_\rho^{\sA\sB}e^{i\pi\ketbra{0}{0}^{\sA\sB\sD}}(\til{W}_\tF^{\sB\sD}U_\rho^{\sA\sB})^\dag U_\sigma^{\sA\sB}e^{i\pi\ketbra{0}{0}^{\sA\sB\sD}}(U_\sigma^{\sA\sB})^\dag,
\end{align}
with the input state $\til{W}_\tF^{\sB\sD}U_\rho^{\sA\sB} \ket{0}^{\sA\sB\sD}$, we can estimate $\sqrt{\rF'}$ within the error $\delta$ using $\cO(1/\delta)$ applications of $\til{W}_\tF^{\sB\sD}U_\rho^{\sA\sB}$, $U_\sigma^{\sA\sB}$, and their inverses.
That is, the output $\sqrt{\til{\rF}} = \mathtt{SqrtAmpEstQuery}(\til{W}_\tF^{\sB\sD}U_\rho^{\sA\sB}, U_\sigma^{\sA\sB}; \delta/2)$ satisfies that 
\begin{align}
\label{inteq:54}
    \Big|\sqrt{\til{\rF}} - \sqrt{\rF'}\Big| \leq \delta/2,
\end{align}
with probability at least $2/3$.


Therefore, we conclude that $\Big|\sqrt{\til{\rF}} - \sqrt{\rF}\Big| \leq \delta$, where we used the triangle inequality and Eqs.~\eqref{inteq:53} and~\eqref{inteq:54}.
The number of query $q$ to $U_\rho^{\sA\sB}$, $U_\sigma^{\sA\sB}$, and their inverses is determined as
\begin{align}
\label{inteq:58}
    q = \cO\big(\log{(1/\delta)}/\beta_\tF\big) \cO(1/\delta) =
    \cO\Big(\f{1}{\delta}{\rm min}\Big\{\f{1}{s_{\rm min}}, \f{r}{\delta}\Big\}\log{\Big(\f{1}{\delta}\Big)}\Big).
\end{align}
In the quantum circuit of this algorithm, the unitary $Q^{\sA\sB\sD}$ in Eq.~\eqref{inteq:48} is repeated $\cO(1/\delta)$ times, where $Q^{\sA\sB\sD}$ includes the Uhlmann transformation consisting of $\cO(\log{(1/\delta)\log{(d_\sA d_\sB)}/\beta_\tF})$ gates (\Cref{thm:Uhlmann alg purif query model}).
Including additional $\cO\big((\log{(1/\delta)})^2\big)$ gates in the square root amplitude estimation, i.e., the quantum phase estimation (\Cref{lem:square root amp est query}), the total number of one- and two-qubit gates is determined by
\begin{align}
    \cO(1/\delta)\cO(\log{(1/\delta)\log{(d_\sA d_\sB)}/\beta_\tF}) + \cO\big((\log{(1/\delta)})^2\big)
    = \cO\big(q\log{(d_\sA d_\sB)} + (\log{(1/\delta)})^2\big).
\end{align}
The number of qubits used at any one time is given by $\cO(\log{(d_\sA d_\sB)} + \log{(1/\delta)})$,
where the first term comes from the Uhlmann transformation and the second term from the square root amplitude estimation.

We complete a proof of Theorem~\ref{thm:fid est quel}.






\end{proof}





%==============================================================================================================================================

\subsection{In the purified and mixed sample access model}
\label{sec:fidelity estimation purif and mixed sample}


We consider the estimation of the fidelity $\sqrt{\rF} = \sqrt{\rF}(\rho^\sA, \sigma^\sA)$ in the purified and mixed sample access models. In \Cref{sec:fidelity estimation collective}, we discuss algorithms which employ collective operations. In Secs.~\ref{sec:fidelity estimation PSS} and~\ref{sec:fidelity estimation SS}, we construct algorithms with non-collective operations in the purified and mixed sample access models, respectively.












\subsubsection{In the purified and mixed sample access models: collective case}
\label{sec:fidelity estimation collective}



We now provide the square root fidelity estimation algorithms in the purified and mixed sample access models when collective operations over multiple samples of the states are allowed.
Our statement is as follows.



\begin{theorem}[Square root fidelity estimation in the purified and mixed sample access models with collective operations]
\label{thm:fidelity est local sample collective}
Let $\delta \in (0, 1)$. Then, in each of the purified and mixed sample access models, there is a quantum sample algorithm that outputs $\sqrt{\til{\rF}}$ such that $\big|\sqrt{\til{\rF}} - \sqrt{\rF}\big| \leq \delta$ with probability at least $2/3$. The algorithm uses $n = \cO\big(\f{1}{\delta^2}{\rm min}\big\{\f{1}{s_{\rm min}^2}, \f{r^2}{\delta^2}\big\}\big(\log{\big(\f{1}{\delta}\big)}\big)^2\big)$ samples of $\ket{\rho}^{\sA\sB}$ and $\ket{\sigma}^{\sA\sB}$ in the purified sample access model, or $\rho^\sA$ and $\sigma^\sA$ in the mixed sample access model.
The quantum circuit of this algorithm consists of $\til{\cO}(n^4)$ one- and two-qubit gates, and $\til{\cO}(n^2)$ qubits suffice at any one time.

\end{theorem}





The key idea in the construction of the algorithms in \Cref{thm:fidelity est local sample collective} is to use a technique recently described in Ref.~\cite{tang2025conjugatequerieshelp}, which aims to simulate a query algorithm by a sample algorithm. This is achieved by using sampled states to approximately implement unitaries that prepare their purified states.
If the original query algorithm attains its goal with probability at least $1-\eta$ using $q$ queries, the corresponding sample algorithm can simulate the same goal with probability at least $1-\eta-\epsilon$ using $\cO(q^2/\epsilon)$ samples.
With this technique, all $q = \cO\big(\f{1}{\delta}{\rm min}\big\{\f{1}{s_{\rm min}}, \f{r}{\delta}\big\}\log{\big(\f{1}{\delta}\big)}\big)$ queries to the purified state preparation unitaries $U_\rho^{\sA\sB}$, $U_\sigma^{\sA\sB}$, and their inverses in the query algorithm $\mathtt{UhlFidelityEstPurifQuery}$ (\Cref{thm:fid est quel}) are replaced by approximate operations constructed from multiple copies of $\ket{\rho}^{\sA\sB}$ and $\ket{\sigma}^{\sA\sB}$, or $\rho^\sA$ and $\sigma^\sA$. Consequently, we then obtain the result in \Cref{thm:fidelity est local sample collective}.


We formally introduce the technique in Ref.~\cite{tang2025conjugatequerieshelp}, slightly adapted for our purposes.
In the following, we focus on the discussion of the mixed sample access model.
For the purified sample access model, the step of random purification described below can simply be omitted.




\begin{theorem}[Simulating queries to state preparation unitary given copies of the state~{\cite[Theorem 1.5]{tang2025conjugatequerieshelp}}]
\label{thm:sample to query lift}
Let $\epsilon, \eta \in (0, 1)$ and let $U_\omega^{\sA\sB}$ be a unitary such that $U_\omega^{\sA\sB}\ket{0}^{\sA\sB} = \ket{\omega}^{\sA\sB}$. Suppose that there is a quantum query algorithm that makes $q$ queries to $U_\omega^{\sA\sB}$ and $(U_\omega^{\sA\sB})^\dag$, and produces an output with probability $p^{\rm qry}$.
If the query algorithm satisfies $p^{\rm qry} \geq 1 - \eta$ for any choice of the purifying system $\sB$, then there exists a quantum sample algorithm that simulates the output of the query algorithm with probability $p^{\rm samp} \geq 1 - \eta - \epsilon$, using $n = \cO(q^2 / \epsilon)$ samples of the state $\omega^\sA$.
The quantum circuit of the sample algorithm consists of $\til{\cO}(n^4)$ one- and two-qubit gates, and $\til{\cO}(n^2)$ qubits at any one time, in addition to those of the query algorithm.



\end{theorem}



At the heart of this algorithm are the following three primitives.
The first primitive is called random purification, which is regarded as an algorithmic strengthening of the results in Refs.~\cite{Soleimanifar2022testingmatrixproduct, chen2024localtestUniInv}. 
The term \emph{random} refers to the following: let $\ket{\omega_0}^{\sA\sB}$ be a fixed purified state of $\omega^\sA$. Then, all purified states of $\omega^\sA$ on $\sA\sB$ are connected by a unitary acting on the purifying system $\sB$; that is, for any purified state, there is a unitary $U^\sB$ such that it can be written as $U^\sB\ket{\omega_0}^{\sA\sB}$.
The random means that the unitary $U^\sB$ is drawn uniformly at random according to the Haar measure. 
The random purification subroutine then allows us to approximately map $(\omega^\sA)^{\otimes n}$ to a state $\bE_U[(U^\sB\ketbra{\omega_0}{\omega_0}^{\sA\sB}(U^{\sB})^\dag)^{\otimes n}]$, i.e., to a state uniformly averaged over the choice of purification.
Remarkably, although the state is averaged, the algorithm prepares $n$ purified states from $n$ samples of $\omega^\sA$ without reducing the number of samples.
The primary component of this subroutine is the Schur transform~\cite{harrow2005applicationSchurtrans, Bacon2006efficientcircuitSchur, Krovi2019efficienthigh, Nguyen2024theorynuralnet, burchardt2025highdimensionalschur}, which mainly determines the number of additional gates and qubits to the circuit of the query algorithm as $\til{\cO}(n^4)$ and $\til{\cO}(n^2)$~\cite{burchardt2025highdimensionalschur}, respectively.


The second primitive, which appeared earlier in Refs.~\cite{kretschmer2021pseudorandomnessclassical, Goldin2025translatingcommonhaar}, relates to a state with a phase: $\ket{\omega_\theta}^{\sA\sB\sS} = \f{1}{\sqrt{2}}(e^{i\theta}\ket{0}^{\sA\sB}\ket{0}^{\sS} + \ket{\omega}^{\sA\sB}\ket{1}^{\sS})$, where $\sS$ is a one-qubit auxiliary system. 
Using this primitive, we can map $(\ketbra{\omega}{\omega}^{\sA\sB})^{\otimes n}$ to $\bE_\theta[(\ketbra{\omega_\theta}{\omega_\theta}^{\sA\sB\sS})^{\otimes n}]$, where the expectation $\bE_\theta[\cdot]$ is taken over $\theta$ chosen uniformly at random from $[0, 2\pi)$. 
The quantum circuit implementing this primitive consists of $\cO(n^2\log{(d_\sA d_\sB)})$ one- and two-qubit gates and uses $\cO(n\log{(d_\sA d_\sB)})$ qubits at any one time \cite{tang2025conjugatequerieshelp}.


The last primitive is the density matrix exponentiation (\Cref{prevthm:cntrl-DME} with $t=\pi$), which maps $(\ketbra{\omega_\theta}{\omega_\theta}^{\sA\sB\sS})^{\otimes m}$ to $e^{i\pi \ketbra{\omega_\theta}{\omega_\theta}^{\sA\sB\sS}}$ within a diamond norm error $\epsilon'$, where $m = \cO(1/\epsilon')$. 
Note that the unitary $e^{i\pi \ketbra{\omega_\theta}{\omega_\theta}^{\sA\sB\sS}}$ is a purified state preparation unitary because $e^{i\pi \ketbra{\omega_\theta}{\omega_\theta}^{\sA\sB\sS}} \ket{0}^{\sA\sB\sS} = -e^{-i\theta} \ket{\omega}^{\sA\sB} \ket{1}^\sS$. Here, the state $-e^{-i\theta} \ket{\omega}^{\sA\sB} \ket{1}^\sS$ is one of the purified states of $\omega^\sA$ with purifying system $\sB\sS$ of dimension at least twice the rank of $\omega^\sA$.
To apply the unitary for $q$ times with total error at most $\epsilon$, we set $\epsilon' = \epsilon/q$ and use $n = qm = \cO(q^2/\epsilon)$ copies of $\ket{\omega_\theta}^{\sA\sB\sS}$.
Naturally, the inverse of the purified state preparation unitary is obtained in a similar manner.


Now, consider a query algorithm that outputs a certain result with probability $p^{\rm qry}$, using $q$ queries to a unitary preparing a purified state of $\omega^\sA$.
We suppose that this query algorithm is independent of the choice of purification, i.e., it satisfies $p^{\rm qry} \geq 1 - \eta$ for any given purified state preparation unitary.
Without loss of generality, we then consider a sufficiently large purifying system, at least twice the rank of $\omega^\sA$.
By replacing all $q$ queries in the query algorithm with the three primitives above, we obtain a sample algorithm.
Using $n = \cO(q^2/\epsilon)$ copies of the state $\omega^\sA$, the sample algorithm can simulate the original query algorithm with probability $p^{\rm samp}$, such that $\big|p^{\rm samp} - \bE_\theta\bE_U [p^{\rm qry}]\big| \leq \epsilon$.
Since $p^{\rm qry} \geq 1 - \eta$ for any $U^\sB$ and $\theta$, it holds that $p^{\rm samp} \geq \bE_\theta \bE_U [p^{\rm qry}] - \epsilon \ge 1 - \eta - \epsilon$.
Hence, \Cref{thm:sample to query lift} follows.





As mentioned in Ref.~\cite{tang2025conjugatequerieshelp}, we should note a caveat: the query algorithm to which we apply \Cref{thm:sample to query lift} must work on any choice of purified state preparation unitary. Since the first and second primitives involve the averages $\bE_\theta[\cdot]$ and $\bE_U[\cdot]$ over possible purifications, the result cannot depend on a specific purification. 
While this condition is typically satisfied in common query algorithms, depending on tasks, a particular purification may be more advantageous than others. For instance, this is the case for the Uhlmann transformation itself. Specifically, the Uhlmann transformation works only for the specific purifications determined by the given initial and final states, $\ket{\rho}^{\sA\sB}$ and $\ket{\sigma}^{\sA\sB}$, and not for arbitrary purifications. On the other hand, when the goal is to estimate the fidelity, this issue can be circumvented by applying the purified state preparation unitaries $U_\rho^{\sA\sB}$ and $U_\sigma^{\sA\sB}$ for the initial and final states, along with the Uhlmann transformation (see also Eqs.~\eqref{inteq:170} to~\eqref{inteq:48}). For this reason, the sample cost of fidelity estimation in the mixed sample access model can be improved by this approach, whereas that of the Uhlmann transformation itself cannot.





Given the above results, \Cref{thm:sample to query lift}, we can straightforwardly prove \Cref{thm:fidelity est local sample collective}.
Note that, unlike other algorithms discussed in this paper, the algorithm in \Cref{thm:sample to query lift}---and thus our algorithms in \Cref{thm:fidelity est local sample collective}---need to handle $n$ samples of the state collectively.







\begin{proof}[Proof of Theorem~\ref{thm:fidelity est local sample collective}]
    

Recalling that for $\delta \in (0, 1)$, and $\eta \in (0, 1)$ of constant order, the algorithm
\begin{align}
    \mathtt{UhlFidelityEstPurifQuery}(U_\rho^{\sA\sB}, U_\sigma^{\sA\sB}; \delta)
\end{align}
in \Cref{alg:fidelestquery} outputs an estimate $\sqrt{\til{\rF}}$ that satisfies $\big|\sqrt{\til{\rF}} - \sqrt{\rF}\big| \leq \delta$ with success probability $p^{\rm qry} \ge 1-\eta$, using $q = \cO\big(\f{1}{\delta}{\rm min}\big\{\f{1}{s_{\rm min}}, \f{r}{\delta}\big\}\log{\big(\f{1}{\delta}\big)}\big)$ queries to $U_\rho^{\sA\sB}$, $U_\sigma^{\sA\sB}$, and their inverses.
Although the success probability was set to at least $2/3$ in \Cref{thm:fid est quel}, it can be amplified to $1-\eta$ with only a constant increase of the query cost as long as $\eta$ is of constant order (see Sec.~\ref{sec:fidelity estimation PQA}).
Note that the query algorithm $\mathtt{UhlFidelityEstPurifQuery}$ works on \emph{any} choice of purification, and hence, outputs the desired estimate with probability at least $1 - \eta$ independent of which purification is selected.


Following Theorem~\ref{thm:sample to query lift}, we implement the query algorithm with $q$ queries by the corresponding sample algorithm by approximately constructing the purified state preparation unitaries.
Using $n = \cO(q^2/\epsilon)$ samples of the states $\rho^\sA$ and $\sigma^\sA$, this sample algorithm outputs $\sqrt{\til{\rF}}$ satisfying $\big|\sqrt{\til{\rF}} - \sqrt{\rF}\big| \leq \delta$ with probability $p^{\rm samp}$ such that $p^{\rm samp} \geq 1 - \eta - \epsilon$.
By setting $\eta = 1/6$ and $\epsilon = 1/6$, we conclude that, for any $\delta \in (0, 1)$, the sample algorithm outputs $\sqrt{\til{\rF}}$ that satisfies $\big|\sqrt{\til{\rF}} - \sqrt{\rF}\big| \leq \delta$ with probability $p^{\rm samp} \geq 2/3$, using $n$ samples of the states $\rho^\sA$ and $\sigma^\sA$, where
\begin{align}
    n 
    &= \cO(q^2/\epsilon) \\
    &=\Big(\f{1}{\delta^2}{\rm min}\Big\{\f{1}{s_{\rm min}^2}, \f{r^2}{\delta^2}\Big\}\Big(\log{\Big(\f{1}{\delta}\Big)}\Big)^2\Big).
\end{align}




We now evaluate the number of gates and qubits in the quantum circuit of this sample algorithm.
The original query algorithm uses $\cO\big(q\log(d_\sA d_\sB) + (\log(1/\delta))^2\big)$ gates and $\cO(\log(d_\sA d_\sB) + \log(1/\delta))$ qubits at once. (\Cref{thm:fid est quel}). 
In addition, $n\log(d_\sA d_\sB)$ qubits are required to store the $n$ samples, so the number of qubits becomes $n\log(d_\sA d_\sB) + \cO(\log(d_\sA d_\sB) + \log(1/\delta)) = \cO(n\log(d_\sA d_\sB))$.
From \Cref{thm:sample to query lift}, converting the query algorithm to the sample one takes $\til{\cO}(n^4)$ gates and $\til{\cO}(n^2)$ qubits simultaneously, in the part of the Schur transform~\cite{burchardt2025highdimensionalschur}.
Therefore, the total number of one- and two-qubit gates in the whole sample algorithm is determined by their sum: $\cO\big(q\log(d_\sA d_\sB) + (\log(1/\delta))^2\big) + \til{\cO}(n^4) = \til{\cO}(n^4)$.
The maximum number of qubits used at one time is given by $\max\{\cO(n\log(d_\sA d_\sB)), \til{\cO}(n^2)\} = \til{\cO}(n^2)$.





\end{proof}








\subsubsection{In the purified sample access model: non-collective case}
\label{sec:fidelity estimation PSS}


Given multiple copies of $\ket{\rho}^{\sA\sB}$ and $\ket{\sigma}^{\sA\sB}$, we aim to estimate the fidelity $\sqrt{\rF}=\sqrt{\rF}(\rho^\sA, \sigma^\sA)$ without using collective operations.
We propose a quantum sample algorithm $\mathtt{UhlFidelityEstPurifSample}$, for which the following theorem guarantees its performance.
The algorithm is described in \Cref{alg:fidelestsample}, where $\sF=\sH\hat{\sA}\hat{\sB}$.


\begin{theorem}[Square root fidelity estimation in the purified sample access model without collective operations]
\label{thm:sqrtfidelestpufifsample}
Let $\delta \in (0, 1)$. Then, the quantum sample algorithm $\mathtt{UhlFidelityEstPurifSample}(\ket{\rho}^{\sA\sB}, \ket{\sigma}^{\sA\sB}; \delta)$,
which is given by \Cref{alg:fidelestsample}, outputs $\sqrt{\til{\rF}}$ such that $\big|\sqrt{\til{\rF}} - \sqrt{\rF}\big| \leq \delta$ with probability at least $2/3$, using $n = \cO\big(\f{1}{\delta^3}{\rm min}\big\{\f{1}{s_{\rm min}^2}, \f{r^2}{\delta^2}\big\}\big(\log{\big(\f{1}{\delta}\big)}\big)^2\big)$ samples of $\ket{\rho}^{\sA\sB}$ and $\ket{\sigma}^{\sA\sB}$.
The quantum circuit of this algorithm consists of $\cO(n\log{(d_\sA d_\sB)})$ one- and two-qubit gates, and $\cO(\log{(d_\sA d_\sB)} + \log{(1/\delta)})$ qubits suffice at any one time.
\end{theorem}


 





\begin{algorithm}[h]
\caption{Square root fidelity estimation algorithm in the purified sample access model \\ \parbox{\linewidth}{\centering $\mathtt{UhlFidelityEstPurifSample}(\ket{\rho}^{\sA\sB}, \ket{\sigma}^{\sA\sB}; \delta)$ (In Theorem~\ref{thm:sqrtfidelestpufifsample})}}
\label{alg:fidelestsample}
\SetKwInput{KwInput}{Input}
\SetKwInput{KwOutput}{Output}
\SetKwInput{KwParameters}{Parameters}
\SetKwComment{Comment}{$\triangleright$\ }{}
\SetCommentSty{textnormal}
\SetKwProg{Fn}{Subroutine}{}{end}
\SetKwFunction{UPS}{UhlmannPurifiedSample}
\SetKwFunction{SAES}{SqrtAmpEstSample}

\SetAlgoNoEnd
\SetAlgoNoLine
\KwInput{Two pure quantum state $\ket{\rho}^{\sA\sB}$ and $\ket{\sigma}^{\sA\sB}$.}
\KwParameters{$\delta \in (0, 1)$.}
\KwOutput{Real number $\sqrt{\til{\rF}}$.}
\SetAlgoLined

Set $\delta_2 \gets \delta/2$ and $\delta_1 \gets \delta_2/(120\pi)$. \\
Set $\cJ_\tF^{\sB\sF} \gets \mathtt{UhlmannPurifSample}(\ket{\rho}^{\sA\sB}, \ket{\sigma}^{\sA\sB}; \delta_1, \tF)$ (Algorithm~\ref{alg:Uhl purif sample}). \\
Set $\omega^{\sA\sB\sF} \gets \cJ_\tF^{\sB\sF}\circ\cP_{\ket{0}\ket{\sigma}}^{\bC\rarr\sF}(\ketbra{\rho}{\rho}^{\sA\sB})$ and $\ket{\psi}^{\sA\sB\sF} \gets \ket{\sigma}^{\sA\sB}\ket{0}^\sH\ket{\rho}^{\hat{\sA}\hat{\sB}}$. \\
Set $\sqrt{\til{\rF}} \gets \mathtt{SqrtAmpEstSample}(\omega^{\sA\sB\sF}, \ket{\psi}^{\sA\sB\sF}; \delta_2)$ (Lemma~\ref{lem:modif square root est sample}). \\
Return $\sqrt{\til{\rF}}$.
\end{algorithm}



The estimation strategy is basically the same as in the last section.
The difference is that, unlike in the purified query access model, we cannot exactly prepare the target unitary $Q$, which is a product of reflections, for the quantum phase estimation.
In this case, a sample version of $\mathtt{SqrtAmpEstQuery}$, i.e., the square root amplitude estimation algorithm for pure state samples $\mathtt{SqrtAmpEstSample}$~\cite{wang2024sampleoptimal}, is helpful.

\begin{theorem}[Square root amplitude estimation algorithm for pure state samples~{\cite[Theorem 7.1]{wang2024sampleoptimal}}]
\label{lem:square root est sample}
For two pure states $\ket{\psi}$ and $\ket{\phi}$ in a $d$-dimensional Hilbert space, and any $\delta \in (0, 1/2)$, there exists a quantum sample algorithm $\mathtt{SqrtAmpEstSample}(\ket{\phi}, \ket{\psi}; \delta)$ that outputs $\sqrt{\til{c}}$ such that $\big|\sqrt{\til{c}} - |\braket{\phi}{\psi}|\big| \leq \delta$ holds with probability at least $2/3$, using $m = \cO\big(1/\delta^2\big)$ samples of $\ket{\psi}$ and $\ket{\phi}$.
The quantum circuit of this algorithm consists of $\cO(m\log{d})$ one- and two-qubit gates, and $\cO(\log{d} + \log{(1/\delta)})$ qubits suffice at any one time.
\end{theorem}

The key idea of this algorithm is as follows. As mentioned in the last section, the quantum phase estimation on the product of reflections $Q = e^{i\pi\ketbra{\psi}{\psi}}e^{i\pi\ketbra{\phi}{\phi}}$ with the input $\ket{\psi}$ and a simple post-processing calculation yield a good estimate of $|\braket{\phi}{\psi}|$.
Since we cannot exactly implement the unitaries $e^{i\pi\ketbra{\psi}{\psi}}$ and $e^{i\pi\ketbra{\phi}{\phi}}$ in the sample model, we instead approximate them with error $\epsilon$ via the density matrix exponentiation, using $\cO(1/\epsilon)$ samples of $\ket{\psi}$ and $\ket{\phi}$. 


Typically, when $\ket{\psi}$ is an eigenstate of $Q$, the quantum phase estimation applies $Q$ a total of $g = \sum_{j=1}^l 2^{j-1} = 2^l - 1 = \cO(1/(\eta\delta))$ times, and uses $\cO\big((\log{(1/\delta)})^2\big)$ gates and $\cO(\log{(1/\delta)})$ auxiliary qubits, to estimate the corresponding phase of $Q$ with error at most $\delta$ and success probability at least $1 - \eta$, where $l = \big\lceil\log{\big(\f{1}{\delta}(2+\f{1}{2\eta})\big)}\big\rceil$~\cite{kitaev1995quantummeasurementsabelian, cleve1997quantumalgorithmrevisit, nielsen2010quantum}.
However, if $Q$ is only approximately prepared, the accumulated error reduces the success probability to $p_{\rm succ} \geq 1 - \eta - 2g\epsilon$.
To ensure that $p_{\rm succ}$ remains a constant order, such as $2/3$, it suffices to set $\eta = 1/6$ and $\epsilon \leq 1/(12g)$, resulting in a total sample complexity of $\ket{\psi}$ and $\ket{\phi}$ as $g\cO(1/\epsilon) = \cO(1/\delta^2)$.
Since the quantum phase estimation requires a controlled version of the target unitary $Q$, we can simply use the controlled implementation of the density matrix exponentiation~\cite{kimmel2017HamSimSampleComp}.


We should note that, in the original paper~\cite{wang2024sampleoptimal}, both inputs to $\mathtt{SqrtAmpEstSample}$ are limited to pure states, while, in \Cref{alg:fidelestsample}, one of them is rather a mixed state $\omega^{\sA\sB\sF}$, which may be close to pure.
This implies that every use of the pure state in $\mathtt{SqrtAmpEstSample}$ is replaced with $\omega^{\sA\sB\sF}$.
We need to evaluate how the error propagates under this modification.

\begin{lemma}[Modified version of Theorem~\ref{lem:square root est sample}]
\label{lem:modif square root est sample}
Suppose that a quantum state $\omega$ satisfies $\f{1}{2}\big\|\omega - \ketbra{\phi}{\phi}\big\|_1 \leq \epsilon'$. Then, for $\delta \in (120\pi\epsilon', 1/2)$ and $\epsilon' \in (0, 1/(240\pi))$, the quantum sample algorithm 
\begin{align}
    \mathtt{SqrtAmpEstSample}(\omega, \ket{\psi}; \delta)
\end{align}
outputs $\sqrt{\til{c}}$ such that $\big|\sqrt{\til{c}} - |\braket{\phi}{\psi}|\big| \leq \delta$ with probability at least $2/3$, using $m = \cO\big(1/\big(\delta(\delta/120 - \pi\epsilon')\big)\big)$ samples of $\omega$ and $\ket{\psi}$.
The quantum circuit of this algorithm consists of $\cO(m\log{d})$ one- and two-qubit gates, and $\cO(\log{d} + \log{(1/\delta)})$ qubits suffice at any one time.
\end{lemma}


\begin{proof}[Proof of \Cref{lem:modif square root est sample}]
We observe that 
\begin{align}
    \f{1}{2}\|e^{i\pi\omega}(\cdot)e^{-i\pi\omega} - e^{i\pi\ketbra{\phi}{\phi}}(\cdot)e^{-i\pi\ketbra{\phi}{\phi}}\|_\diamond
    &\leq \big\|e^{i\pi\omega} - e^{i\pi\ketbra{\phi}{\phi}}\big\|_\infty \\
    &\leq \pi \big\|\omega - \ketbra{\phi}{\phi}\big\|_\infty \\
    &\leq \pi \big\|\omega - \ketbra{\phi}{\phi}\big\|_1 \\
    &\leq 2\pi \epsilon',
\end{align}
where we used \Cref{lem:state error not accumulate} in the second inequality, and the last inequality follows from the assumption.
Using the density matrix exponentiation with $\cO(1/\epsilon)$ samples of $\omega$, we can implement the quantum channel $\cL$ that satisfies $\f{1}{2}\big\|\cL - e^{i\pi\omega}(\cdot)e^{-i\omega}\big\|_\diamond \leq \epsilon$.
Thus, we have that
\begin{align}
    \f{1}{2}\big\|\cL - e^{i\pi\ketbra{\phi}{\phi}}(\cdot)e^{-i\pi\ketbra{\phi}{\phi}}\big\|_\diamond 
    &\leq \f{1}{2}\|\cL - e^{i\pi\omega}(\cdot)e^{-i\pi\omega}\|_\diamond  \notag\\
    &\hspace{2pc}+\f{1}{2}\|e^{i\pi\omega}(\cdot)e^{-i\pi\omega} - e^{i\pi\ketbra{\phi}{\phi}}(\cdot)e^{-i\pi\ketbra{\phi}{\phi}}\|_\diamond \\
    &\leq \epsilon + 2\pi \epsilon'.
\end{align}

We recall that in the quantum phase estimation on the product of reflections $Q = e^{i\pi\ketbra{\psi}{\psi}} e^{i\pi\ketbra{\phi}{\phi}}$, the target unitary $Q$ is used $g = 2^l -1 $ times, where $l = \big\lceil\log{\big(\f{1}{\delta}(2+\f{1}{2\eta})\big)}\big\rceil$.
When we replace all $g$ uses of $e^{i\pi\ketbra{\phi}{\phi}}$ with the channel $\cL$, the resulting error is $g(\epsilon + 2\pi \epsilon')$.
Combining this with the error $g\epsilon$ from approximating $g$ uses of $e^{i\pi\ketbra{\psi}{\psi}}$ via the density matrix exponentiation with $\cO(1/\epsilon)$ samples of $\ket{\psi}$, a total error becomes $2g(\epsilon + \pi \epsilon')$. 

Hence, the quantum phase estimation outputs $\sqrt{\til{c}}$ that satisfies $\big|\sqrt{\til{c}} - |\braket{\psi}{\phi}|\big| \leq \delta$ with probability $p_{\rm succ} \geq 1 - \eta - 2g(\epsilon + \pi \epsilon')$.
Setting $\eta = 1/6$ and $\epsilon = 1/(12g) - \pi\epsilon'$, the success probability is at least $2/3$.
To ensure that $\epsilon > 0$, it suffices to take $\delta > 120\pi\epsilon'$, since $g = 2^{\lceil \log{(5/\delta)}\rceil}-1 < 10/\delta$.
The algorithm uses $g\cO(1/\epsilon) = \cO\Big(1/\big(\delta(\delta/120 - \pi\epsilon')\big)\Big)$ samples of $\omega$ and $\ket{\psi}$.
\end{proof}



With the above discussion, we now prove Theorem~\ref{thm:sqrtfidelestpufifsample}.



\begin{proof}[Proof of Theorem~\ref{thm:sqrtfidelestpufifsample}]



Let $\sF=\sH\hat{\sA}\hat{\sB}$.
Through the Uhlmann transformation algorithm, we obtain a quantum channel $\cJ_\tF^{\sB\sF} = \mathtt{UhlmannPurifiedSample}(\ket{\rho}^{\sA\sB}, \ket{\sigma}^{\sA\sB}; \delta_1, \tF)$, using $\cO\big(\big(\log{(1/\delta_1)}\big)^2/(\delta_1\beta_\tF^2)\big)$ samples.
We here chose $\beta_\tF$, which is given by $\beta_\tF = \f{2}{\pi}\max\Big\{s_{\rm min}, \f{2\delta_1}{r}\Big\}$. Recall that $s_{\rm min}$ and $r$ are the minimum non-zero singular value and the rank of $\sqrt{\sigma^\sA}\sqrt{\rho^\sA}$, respectively.


As discussed in \Cref{sec:purif sample}, the channel $\cJ_\tF^{\sB\sF}$ satisfies that
\begin{align}
\label{inteq:69}
    \f{1}{2}\big\|\cJ_\tF^{\sB\sF} - \cU_1^{\sB\sF}\big\|_\diamond \leq \delta_1/2,
\end{align}
where $U_1^{\sB\sF}$ is a block-encoding unitary of $\ketbra{\rho}{\sigma}^{\hat{\sA}\hat{\sB}}\otimes P_\sgn^{(\rm SV)}\big(\sin^{(\rm SV)}(M^\sB)\big)$.
(See Eq.~\eqref{inteq:16}. Note that in $\mathtt{UhlmannPurifiedSample}$ the parameter has already been rescaled as $\delta_2 \gets \delta_1 / (2u)$.)  
As explained in \Cref{sec:uhlfidpurifsamp}, when we consider the difference between $\sqrt{\rF}=\sqrt{\rF}(\rho^\sA, \sigma^\sA)$ and 
\begin{align}
    \sqrt{\rF'} = \sqrt{\rF}(U_1^{\sB\sF}\ket{\rho}^{\sA\sB}\ket{0}^\sH\ket{\sigma}^{\hat{\sA}\hat{\sB}}, \ket{\sigma}^{\sA\sB}\ket{0}^\sH\ket{\rho}^{\hat{\sA}\hat{\sB}}),
\end{align}
we see that $\big|\sqrt{\rF} - \sqrt{\rF'}| \leq \delta_1/4$ (see Eq.~\eqref{inteq:51} with $\delta_1$ rescaled as $\delta_1 \gets \delta_1/8$). 


Next, we consider estimating $\sqrt{\rF'}$. When we define $\omega^{\sA\sB\sF} = \cJ_\tF^{\sB\sF}\circ\cP_{\ket{0}\ket{\sigma}}^{\bC\rarr\sF}(\ketbra{\rho}{\rho}^{\sA\sB})$ and $\ket{\phi}^{\sA\sB\sF} = U_1^{\sB\sH}\ket{\rho}^{\sA\sB}\ket{0}^\sH\ket{\sigma}^{\hat{\sA}\hat{\sB}}$, Eq.~\eqref{inteq:69} implies that 
\begin{equation}
\label{eq:state error in FE}
    \f{1}{2}\big\|\omega^{\sA\sB\sF} - \ketbra{\phi}{\phi}^{\sA\sB\sF}\big\|_1 \leq \delta_1/2. 
\end{equation} 
Let $\sqrt{\til{\rF}} = \mathtt{SqrtAmpEstSample}(\omega^{\sA\sB\sF}, \ket{\psi}^{\sA\sB\sF}; \delta_2)$ for $\delta_2 \in (60\pi \delta_1, 1/2)$, where $\ket{\psi}^{\sA\sB\sF} = \ket{\sigma}^{\sA\sB}\ket{0}^\sH\ket{\rho}^{\hat{\sA}\hat{\sB}}$.
From \Cref{lem:modif square root est sample}, we observe that $\big|\sqrt{\til{\rF}} - \sqrt{\rF'}\big| \leq \delta_2$ with probability at least $2/3$, using $\cO\Big(1/\big(\delta_2(\delta_2/120 - \pi \delta_1/2)\big)\Big)$ samples of $\omega^{\sA\sB\sF}$ and $\ket{\phi}^{\sA\sB\sF}$.
Using the triangle inequality, we can obtain that
\begin{align}
\label{inteq:72}
    \Big|\sqrt{\til{\rF}} - \sqrt{\rF}\Big| 
    &\leq \Big|\sqrt{\til{\rF}} - \sqrt{\rF'}\Big| + \big|\sqrt{\rF'} - \sqrt{\rF}\big| \\
    &\leq \delta_2 + \delta_1/4.
\end{align}
By choosing $\delta_1 = \delta_2/(120\pi)$ and $\delta_2 = \delta/2$, we have 
\begin{align}
    \Big|\sqrt{\til{\rF}} - \sqrt{\rF}\Big| 
    &\leq \f{1 + 480\pi}{960\pi}\delta \\
    &\leq \delta.
\end{align}
We conclude that $\sqrt{\rF}$ is estimated with accuracy $\delta$ with probability at least $2/3$.

The number of samples $n$ of $\ket{\rho}^{\sA\sB}$ and $\ket{\sigma}^{\sA\sB}$ is given by
\begin{align}
    n = \cO\big(\big(\log{(1/\delta_1)}\big)^2/(\delta_1\beta_\tF^2)\big)\cO\big(1/\big(\delta_2(\delta_2/120 - \pi \delta_1/2)\big)\big)   =\cO\Big(\f{1}{\delta^3}\min\Big\{\f{1}{s_{\rm min}^2}, \f{r^2}{\delta^2}\Big\}\Big(\log{\Big(\f{1}{\delta}\Big)}\Big)^2\Big).
\end{align}
In the quantum circuit of this algorithm, the Uhlmann transformation, which consists of $\cO\big((\log{(1/\delta)})^2\allowbreak\log{(d_\sA d_\sB)}/(\delta\beta_\tF^2)\big)$ gates (\Cref{thm:algorithm Uhlmann}) is used $m = \cO(1/\delta^2)$ times, and additional $\cO\big(m\log{(d_\sA d_\sB)}\big)$ gates are used for the square root amplitude estimation (\Cref{lem:modif square root est sample}). Thus, the total number of one- and two-qubit gates is given by $m\cO\big((\log{(1/\delta)})^2\log{(d_\sA d_\sB)}/(\delta\beta_\tF^2)\big) + \cO(m\log{(d_\sA d_\sB)}) = \cO(n\log{(d_\sA d_\sB)})$.
At any one time in this algorithm, $\cO(\log{(d_\sA d_\sB)} + \log{(1/\delta)})$ qubits suffice, considering the Uhlmann transformation and the square root amplitude estimation, i.e., the quantum phase estimation.



We complete a proof of Theorem~\ref{thm:sqrtfidelestpufifsample}.


    
\end{proof}








%=============================================================================================================================================================






\subsubsection{In the mixed sample access model: non-collective case}
\label{sec:fidelity estimation SS}


We consider the estimation of the fidelity $\sqrt{\rF} = \sqrt{\rF}(\rho^\sA, \sigma^\sA)$, where we are given many identical copies of $\rho^\sA$ and $\sigma^\sA$. Collective operations are not allowed here.
The algorithm $\mathtt{UhlFidelityEstMixedSample}$ is provided in \Cref{alg:fidelestsamplemixed}, where $\sF = \sH\hat{\sA}\hat{\sB}$ with $d_\sA = d_\sB$, and $\sH$ is a two-qubit auxiliary system.



\begin{theorem}[Square root fidelity estimation in the mixed sample access model without collective operations]
\label{thm:fidelity est local sample}
Let $\delta \in (0, 1)$. Then, the quantum sample algorithm $\mathtt{UhlFidelityEstMixedSample}(\rho^\sA, \sigma^\sA; \delta)$,
which is given by \Cref{alg:fidelestsamplemixed}, outputs $\sqrt{\til{\rF}}$ such that $\big|\sqrt{\til{\rF}} - \sqrt{\rF}\big| \leq \delta$ with probability at least $2/3$, using $n$ samples of $\rho^\sA$ and $\sigma^\sA$, where
\begin{align}
    n = \til{\cO}\Big(\f{d_\sA}{\delta^4}\min\Big\{\f{\kappa_\rho^2 + \kappa_\sigma^2}{s_{\rm min}^3}, \f{d_\sA^2 r^7}{\delta^{11}}\Big\}\Big).
\end{align}
The quantum circuit of this algorithm consists of $\cO(n\log{d_\sA})$ one- and two-qubit gates, and $\cO(\log{d_\sA} + \log{(1/\delta)})$ qubits suffice at any one time.
\end{theorem}







\begin{algorithm}[h]
\caption{Square root fidelity estimation algorithm in the mixed sample access model \\ \parbox{\linewidth}{\centering $\mathtt{UhlFidelityEstMixedSample}(\rho^{\sA}, \sigma^{\sA}; \delta)$ (In Theorem~\ref{thm:fidelity est local sample})}}
\label{alg:fidelestsamplemixed}
\SetKwInput{KwInput}{Input}
\SetKwInput{KwOutput}{Output}
\SetKwInput{KwParameters}{Parameters}
\SetKwComment{Comment}{$\triangleright$\ }{}
\SetCommentSty{textnormal}
\SetKwProg{Fn}{Subroutine}{}{end}
\SetKwFunction{UPS}{UhlmannPurifiedSample}
\SetKwFunction{SAES}{SqrtAmpEstSample}

\SetAlgoNoEnd
\SetAlgoNoLine
\KwInput{Two quantum states $\rho^{\sA}$ and $\sigma^{\sA}$.}
\KwParameters{$\delta \in (0, 1)$.}
\KwOutput{Real number $\sqrt{\til{\rF}}$.}
\SetAlgoLined

Set $\delta_2 \gets \delta/2$ and $\delta_1 \gets \delta_2/(960\pi)$. \\
Set $\cJ_\tF^{\sB\sF}\circ\cP_{\til{\sigma}_{\rm c}}^{\bC\rarr\hat{\sA}\hat{\sB}} \gets \mathtt{UhlmannMixedSample}(\rho^{\sA}, \sigma^{\sA}; \delta_1, \tF)$ (Algorithm~\ref{alg:Uhl mixed sample}). \\
Set $\til{\rho}_{\rm c}^{\sA\sB} \gets \mathtt{CanonicalPurification}(\rho^\sA; \delta_1/4)$ and  $\til{\sigma}_{\rm c}^{\sA\sB} \gets \mathtt{CanonicalPurification}(\sigma^\sA; 3\delta_1/4)$ (Algorithm~\ref{alg:canonical purification}). \\
Set $\omega^{\sA\sB\sF} \gets \cJ_\tF^{\sB\sF}\circ\cP_{\til{\sigma}_{\rm c}}^{\bC\rarr\hat{\sA}\hat{\sB}}\circ\cP_{\ket{0}}^{\bC\rarr\sH}(\til{\rho}_{\rm c}^{\sA\sB})$ and $\mu^{\sA\sB\sF} \gets \til{\sigma}_{\rm c}^{\sA\sB}\otimes\til{\rho}_{\rm c}^{\hat{\sA}\hat{\sB}}\otimes\ketbra{0}{0}^{\sH}$. \\
Set $\sqrt{\til{\rF}} \gets \mathtt{SqrtAmpEstSample}(\omega^{\sA\sB\sF}, \mu^{\sA\sB\sF}; \delta_2)$ (Lemma~\ref{lem:modif square root est sample}). \\
Return $\sqrt{\til{\rF}}$.


\end{algorithm}



The approach is to first use the canonical purification algorithm $\mathtt{CanonicalPurification}$ and the Uhlmann transformation algorithm in the mixed sample access model $\mathtt{UhlmannMixedSample}$. We then apply the square root amplitude estimation algorithm $\mathtt{SqrtAmpEstSample}$.




\begin{proof}[Proof of Theorem~\ref{thm:fidelity est local sample}]
    



From the discussions in \Cref{sec:mixed sample}, we can obtain a quantum channel
\begin{align}
    \cJ_\tF^{\sB\sF}\circ\cP_{\til{\sigma}_{\rm c}}^{\bC\rarr\hat{\sA}\hat{\sB}} = \mathtt{UhlmannMixedSample}(\rho^\sA, \sigma^\sA; \delta_1, \tF),
\end{align}
which satisfies that $\f{1}{2}\big\|\cJ_\tF^{\sB\sF}\circ\cP_{\til{\sigma}_{\rm c}}^{\bC\rarr\hat{\sA}\hat{\sB}} - \cU_1^{\sB\sF}\circ\cP_{\ket{\sigma_{\rm c}}}^{\bC\rarr\hat{\sA}\hat{\sB}}\big\|_\diamond \leq 3\delta_1/4$, using 
\begin{align}
\label{inteq:75}
    \zeta &= \til{\cO}\Big(\f{d_\sA}{\delta_1^2\beta_\tF^3}\min\Big\{\kappa_\rho^2 + \kappa_\sigma^2, \f{d_\sA^2}{\delta_1^4\beta_\tF^4}\Big\}\Big),
\end{align}
samples of $\rho^\sA$ and $\sigma^\sA$. (See Eq.~\eqref{inteq:100}. Note that in $\mathtt{UhlmannMixedSample}$ the parameters have already been rescaled as $\delta_1 \gets \delta_1/\big(2(4u+1)\big)$ and $\delta_2 \gets \delta_1/(4u)$.)
Here, $U_1^{\sB\sF}$ is a block-encoding unitary of $\ketbra{\rho_{\rm c}}{\sigma_{\rm c}}^{\hat{\sA}\hat{\sB}} \otimes P_\sgn(\tr_{\sA}\big[\ketbra{\sigma_{\rm c}}{\rho_{\rm c}}^{\sA\sB}\big])$, and $\beta_\tF = \cO\big(\max\{s_{\rm min}, \delta_1/r\}\big)$. Moreover, when we define $\sqrt{\rF'}$ by $\sqrt{\rF'} = \sqrt{\rF}(U_1^{\sB\sF}\ket{\rho_{\rm c}}^{\sA\sB}\ket{\sigma_{\rm c}}^{\hat{\sA}\hat{\sB}}\ket{0}^{\sH}, \ket{\sigma_{\rm c}}^{\sA\sB}\ket{\rho_{\rm c}}^{\hat{\sA}\hat{\sB}}\ket{0}^{\sH})$, it holds that $\big|\sqrt{\rF} - \sqrt{\rF'}\big| \leq \delta_1/8$. (See Eq.~\eqref{inteq:77} with $\delta_3$ rescaled as $\delta_3 \gets \delta_1/16$.)


Let $\til{\rho}_{\rm c}^{\sA\sB} = \mathtt{CannonicalPurification}(\rho^\sA; \delta_1/4)$.
Then, we denote by $\omega^{\sA\sB\sF}$ and $\ket{\phi}^{\sA\sB\sF}$ the states given by $\omega^{\sA\sB\sF} = \cJ_\tF^{\sB\sF}\circ\cP_{\til{\sigma}_{\rm c}}^{\bC\rarr\hat{\sA}\hat{\sB}}\circ\cP_{\ket{0}}^{\bC\rarr\sH}(\til{\rho}_{\rm c}^{\sA\sB})$,
and $\ket{\phi}^{\sA\sB\sF} = U_1^{\sB\sF}\ket{\rho_{\rm c}}^{\sA\sB}\ket{\sigma_{\rm c}}^{\hat{\sA}\hat{\sB}}\ket{0}^{\sH}$, respectively.
We see that
\begin{align}
    \f{1}{2}\big\|\omega^{\sA\sB\sF} - \ketbra{\phi}{\phi}^{\sA\sB\sF}\big\|_1
    &\leq \f{1}{2}\big\|\cJ_\tF^{\sB\sF}\circ\cP_{\til{\sigma}_{\rm c}}^{\bC\rarr\hat{\sA}\hat{\sB}} - \cU_1^{\sB\sF}\circ\cP_{\ket{\sigma_{\rm c}}}^{\bC\rarr\hat{\sA}\hat{\sB}}\big\|_\diamond + \f{1}{2}\big\|\til{\rho_{\rm c}}^{\sA\sB} - \ketbra{\rho_{\rm c}}{\rho_{\rm c}}^{\sA\sB}\big\|_1 \\
    &\leq 3\delta_1/4 + \delta_1/4 \\
    \label{inteq:73}
    &=\delta_1.
\end{align}


Let $\mu^{\sA\sB\sF} = \til{\sigma}_{\rm c}^{\sA\sB}\otimes\til{\rho}_{\rm c}^{\hat{\sA}\hat{\sB}}\otimes\ketbra{0}{0}^{\sH}$, where $\til{\sigma}_{\rm c}^{\sA\sB} = \mathtt{CannonicalPurification}(\sigma^\sA; 3\delta_1/4)$.
We evaluate the distance between $\mu^{\sA\sB\sF}$ and $\ket{\psi} = \ket{\sigma_{\rm c}}^{\sA\sB}\ket{\rho_{\rm c}}^{\hat{\sA}\hat{\sB}}\ket{0}^{\sH}$ as
\begin{align}
    \f{1}{2}\big\|\mu^{\sA\sB\sF} - \ketbra{\psi}{\psi}^{\sA\sB\sF}\big\|_1
    &\leq \f{1}{2}\big\|\til{\sigma}_{\rm c}^{\sA\sB} - \ketbra{\sigma_{\rm c}}{\sigma_{\rm c}}^{\sA\sB}\big\|_1 + \f{1}{2}\big\|\til{\rho}_{\rm c}^{\hat{\sA}\hat{\sB}} - \ketbra{\rho_{\rm c}}{\rho_{\rm c}}^{\hat{\sA}\hat{\sB}}\big\|_1 \\
    &\leq 3\delta_1/4 + \delta_1/4 \\
    \label{inteq:74}
    &= \delta_1.
\end{align}


We run $\mathtt{SqrtAmpEstSample}$ with $\omega^{\sA\sB\sF}$ and $\mu^{\sA\sB\sF}$ as inputs, instead of $\ket{\phi}^{\sA\sB\sF}$ and $\ket{\psi}^{\sA\sB\sF}$,  where both inputs are mixed states that are close to pure within error $\delta_1$ (see Eqs.~\eqref{inteq:73} and~\eqref{inteq:74}).
Even in this case, the number of samples can be straightforwardly evaluated in a similar way to the previous section.



Using the density matrix exponentiation, we can implement a quantum channel that approximates $e^{i\pi\ketbra{\phi}{\phi}^{\sA\sB\sF}}$ with error at most $\epsilon + 2\pi\delta_1$, using $\cO(1/\epsilon)$ samples of $\omega^{\sA\sB\sF}$.
Similarly, using $\cO(1/\epsilon)$ samples of $\mu^{\sA\sB\sF}$, we can construct a quantum channel that approximates $e^{i\pi\ketbra{\psi}{\psi}^{\sA\sB\sF}}$ with the same error.
Hence, we apply the quantum phase estimation using these channels instead of
$e^{i\pi\ketbra{\phi}{\phi}^{\sA\sB\sF}}$ and $e^{i\pi\ketbra{\psi}{\psi}^{\sA\sB\sF}}$,
for $g = 2^l - 1$ times, where
$l = \big\lceil\log{\big(\f{1}{\delta_2}\big(2 + \f{1}{2\eta}\big)}\big)\big\rceil$.
This allows us to estimate $|\braket{\psi}{\phi}|$ within the error $\delta_2$ with success probability at least
$1 - \eta - 2g(\epsilon + 2\pi\delta_1) - \delta_1$,
where the last term $-\delta_1$ arises from the fact that the input to the quantum phase estimation is $\mu^{\sA\sB\sF}$ rather than $\ket{\psi}^{\sA\sB\sF}$.

Setting $\eta = 1/6$ and $\epsilon = 1/(24g) - 2\pi \delta_1$, the success probability is at least $3/4 - \delta_1 \geq 2/3$, where we took $\delta_1 \in (0, 1/(960\pi))$, to ensure $\epsilon >0$.
As a result, for $\delta_2 \in (480\pi\delta_1, 1/2)$ and $\delta_1 \in (0, 1/(960\pi))$, $\mathtt{SqrtAmpEstSample}(\omega^{\sA\sB\sF}, \mu^{\sA\sB\sF}; \delta_2)$ outputs $\sqrt{\til{\rF}}$ such that $\big|\sqrt{\til{\rF}} - |\braket{\psi}{\phi}|\big| \leq \delta$ with probability at least $2/3$, using $g\cO(1/\epsilon) = \cO\Big(1/\big(\delta_2(\delta_2/240 - 2\pi \delta_1)\big)\Big)$ samples of $\omega^{\sA\sB\sF}$ and $\mu^{\sA\sB\sF}$.


Noting that $|\braket{\psi}{\phi}| = \sqrt{\rF'}$, we have that
\begin{align}
    \big|\sqrt{\til{\rF}} - \sqrt{\rF}\big| &\leq \big|\sqrt{\rF'} - \sqrt{\rF}\big| + \big|\sqrt{\til{\rF}} - \sqrt{\rF'}\big| \\
    & \leq\delta_1/8 + \delta_2.
\end{align}
Setting $\delta_1 = \delta_2/(960\pi)$ and $\delta_2 = \delta/2$, we obtain that
\begin{align}
    \big|\sqrt{\til{\rF}} - \sqrt{\rF}\big| 
    &\leq \f{1+7680\pi}{15360\pi}\delta \\
    &\leq \delta.
\end{align}
The number of samples $n$ of $\rho^\sA$ and $\sigma^\sA$ is evaluated as
\begin{align}
    n &= \zeta \cO\Big(1/\big(\delta_2(\delta_2/240 - 2\pi \delta_1)\big)\Big) \\
    &=\til{\cO}\Big(\f{d_\sA}{\delta_1^2\beta_\tF^3}\min\Big\{\kappa_\rho^2 + \kappa_\sigma^2, \f{d_\sA^2}{\delta_1^4\beta_\tF^4}\Big\}\Big)\cO\Big(1/\big(\delta_2(\delta_2/240 - 2\pi \delta_1)\big)\Big) \\
    &= \til{\cO}\Big(\f{d_\sA}{\delta^4}
    \min\Big\{\f{\kappa_\rho^2 + \kappa_\sigma^2}{s_{\rm min}^3}, \f{d_\sA^2 r^7}{\delta^{11}}\Big\}\Big).
\end{align}


In the quantum circuit of this algorithm, the most costly component, the Uhlmann transformation, which consists of $\cO\big(\zeta \log(d_\sA d_\sB)/(\delta\beta_\tF^2)\big)$ gates (\Cref{thm:Uhlmann mixed state sample}), is applied $m = \cO(1/\delta^2)$ times. Hence, the total number of one- and two-qubit gates is $m \cO(\zeta \log d_\sA) = \cO(n \log d_\sA)$.
Considering both the Uhlmann transformation and the square root amplitude estimation (i.e., quantum phase estimation), $\cO(\log d_\sA + \log(1/\delta))$ qubits suffice at any one time during this algorithm.



We complete a proof of Theorem~\ref{thm:fidelity est local sample}.


\end{proof}








%================================================================================================================================================





\subsection{Lower bound on the query and sample complexities}
\label{sec:lower bound}


We here provide a lower bound on the query and sample complexities for estimation of the square root fidelity. Note that, since the purified query access model is a more powerful computational model than the purified and mixed sample access models, a lower bound on the query complexity immediately implies corresponding lower bounds on the sample complexities. This lower bound directly leads to a lower bound on the query and sample complexities required to realize the Uhlmann transformation.


Regarding a lower bound on the query complexity for the square root fidelity estimation in the purified query access model, we provide the following proposition.
\begin{proposition}[Lower bound on the query complexity for square root fidelity estimation]
\label{prop:lower fidelity est}
Suppose we are given purified query access to unitaries $U_\rho$, $U_\sigma$, and their inverses.
Then, for $\delta \in (0, 1/16)$, there exists a pair of rank-$k$ states $\rho$ and $\sigma$ such that any quantum query algorithm that estimates $\sqrt{\rF}(\rho, \sigma)$ within error $\delta$ requires a total of $\Omega\big(\max\big\{k^{1/3}, 1/\delta\big\}\big)$ queries to $U_\rho$, $U_\sigma$, and their inverses.
\end{proposition}




The lower bound $\Omega(1/\delta)$ is a direct consequence of the results in Refs.~\cite{belovs2019ClassicalDist, Wang2024OptimalFidelityPure, luo2024SuccinctTest, Liu2024geometricmean}. Note that this also holds for the estimation of fidelity between two pure states, i.e., the case where $k=1$, which implies that the algorithm in \Cref{lem:square root amp est query} is optimal.
To prove \Cref{prop:lower fidelity est}, we use the result in Ref.~\cite{chakraborty2010NewResultQPT}, a lower bound on the query complexity of testing uniformity of a classical probability distribution~\cite{bravi2011PropTest}.
Here, we use the version in Ref.~\cite{Liu2025EstimationTrace}, which applies it to the mixedness testing for a quantum state.





\begin{theorem}[Lower bound for the mixedness testing in the purified query access model~{\cite[Lemma 2.11]{Liu2025EstimationTrace}}]
\label{lem:query lower}
Let $\omega$ be a state of rank-$k$, and let $\pi_k$ be the state with uniform eigenvalues on the support of $\omega$ and zero elsewhere.
Suppose that $U_\omega$ is a unitary that prepares a purified state of $\omega$, i.e., $U_\omega\ket{0} = \ket{\omega}$.
Then, for $\delta \in (0, 1/2]$, any quantum query algorithm that determines whether $\omega = \pi_k$ or $\f{1}{2}\|\omega - \pi_k\|_1 \geq \epsilon$, requires a total of $\Omega(k^{1/3})$ queries to $U_\omega$ and its inverse.
\end{theorem}

Using Theorem~\ref{lem:query lower}, we prove \Cref{prop:lower fidelity est}.

\begin{proof}[Proof of \Cref{prop:lower fidelity est}]

Let $U_\rho^{\sA\sB}$ and $U_{\Phi_k}^{\sA\sB}$ be unitaries such that
\begin{align}
    \ket{\rho}^{\sA\sB} = U_\rho^{\sA\sB}\ket{0}^{\sA\sB} = \sum_{j=1}^k \sqrt{p_j} \ket{j}^\sA\ket{\psi_j}^\sB, \ \ \text{and} \ \ \ 
    \ket{\Phi_k}^{\sA\sB} = U_{\Phi_k}^{\sA\sB}\ket{0}^{\sA\sB} = \frac{1}{\sqrt{k}} \sum_{j=1}^k \ket{j}^\sA\ket{j}^\sB,
\end{align}
where each $\{\ket{j}\}_j$ and $\{\ket{\psi_j}\}_j$ forms an orthonormal basis.
These states are purified states of 
\begin{align}
\label{inteq:164}
    \rho^\sA = \sum_{j=1}^k p_j\ketbra{j}{j}^\sA, \ \ \text{and} \ \ \  \pi_k^\sA = \f{1}{k}\sum_{j=1}^k\ketbra{j}{j}^\sA, 
\end{align}
respectively.

We assume that $ \f{1}{2}\|\rho^\sA - \pi_k^{\sA}\|_1 \geq 1/2$.
By the Fuchs–van de Graaf inequalities, we have $\sqrt{\rF}(\rho^\sA, \pi_k^\sA) \leq \f{\sqrt{3}}{2}$.
Note that $\pi_k^\sA$ is known and $\sqrt{\rF}(\pi_k^\sA, \pi_k^\sA) = 1$.
Hence, any quantum algorithm that can estimate the fidelity $\sqrt{\rF}(\cdot, \pi_k^\sA)$ within error $\delta < \f{1}{16} \ ( < \f{1}{2}(1-\f{\sqrt{3}}{2}))$ can be used to distinguish which unitary of $U_\rho^{\sA\sB}$ and $U_{\Phi_k}^{\sA\sB}$ we applied. 

On the other hand, Theorem~\ref{lem:query lower} states that, to distinguish these unitaries, $\Omega(k^{1/3})$ queries are required, even when $\delta$ is constant.
These imply that there exists a pair of rank-$k$ quantum states such that any fidelity estimation algorithm requires at least $\Omega(k^{1/3})$ queries.


Combining this with the bound of $\Omega(1/\delta)$ yields \Cref{prop:lower fidelity est}.
\end{proof}




We again mention a lower bound on the sample complexity for square root fidelity estimation in the purified and mixed sample access models. 
Since the purified query access model is the most powerful of the three computational models, \Cref{prop:lower fidelity est} implies a corresponding lower bound on the sample complexity.
Hence, there is a pair of rank-$k$ states such that any quantum algorithm for estimating the square root fidelity with error at most $\delta$ requires $\Omega(\max\{k^{1/3}, 1/\delta\})$ samples in the purified and mixed sample access models.
For the mixed sample access model, however, a tighter bound is known~\cite{gilyen2022improvedfidelity}: any quantum algorithm within error $\delta$ requires $\Omega(k/\delta)$ samples for some pair of rank-$k$ states.



As seen in the previous sections, the Uhlmann transformation algorithm can be used to estimate the fidelity. While \Cref{prop:lower fidelity est} is stated for the square root fidelity, its derivation can be straightforwardly adapted to the squared fidelity. 
This leads to \Cref{cor:lower Uhlmann query}, which we restate below.



\corloweruhlquery*



\begin{proof}[Proof of \Cref{cor:lower Uhlmann query}]


Suppose that we can implement a quantum channel $\cT^\sB$ satisfying
\begin{align}
    \label{inteq:152}
    \big|\rF\big(\cT^\sB(\ketbra{\rho}{\rho}^{\sA\sB}), \ket{\sigma}^{\sA\sB}\big) - \rF(\rho^\sA, \sigma^\sA)\big| \leq \delta,
\end{align}
using $q$ queries to $U_\rho^{\sA\sB}$, $U_\sigma^{\sA\sB}$, and their inverses.
Then, by performing the swap test~\cite{Barenco1997stabQCsymm, Buhrman2001fingerprint, Nishimura2025ASurveySWAPtest} on $\cT^\sB(\ketbra{\rho}{\rho}^{\sA\sB})$ and $\ket{\sigma}^{\sA\sB}$ a constant number of times, we can obtain an estimate $\til{\rF}$ such that $\big|\til{\rF} - \rF\big(\cT^\sB(\ketbra{\rho}{\rho}^{\sA\sB}), \ket{\sigma}^{\sA\sB}\big)\big| \leq 1/72$.
Note that since $\ket{\sigma}^{\sA\sB}$ is pure, $\rF\big(\cT^\sB(\ketbra{\rho}{\rho}^{\sA\sB}), \ket{\sigma}^{\sA\sB}\big) = \big|\bra{\sigma}\cT^\sB(\ketbra{\rho}{\rho}^{\sA\sB})\ket{\sigma}^{\sA\sB}\big|$.
Using the triangle inequality, we have $\big|\til{\rF} - \rF(\rho^\sA, \sigma^\sA)\big| \leq \delta + 1/72$.
Hence, we can estimate the fidelity $\rF(\rho^\sA, \sigma^\sA)$ with additive error $\delta + 1/72$, using a constant multiple of $q$ queries.


On the other hand, for certain rank-$k$ states $\rho^\sA$ and $\sigma^\sA$, $\Omega(k^{1/3})$ queries are required to estimate $\rF(\rho^\sA, \sigma^\sA)$ within additive error $1/8$.
This follows by adapting the proof of \Cref{prop:lower fidelity est} to the squared fidelity instead of the square root fidelity. 
In particular, for the states $\rho^\sA$ and $\pi_k^\sA$ given by Eq.~\eqref{inteq:164} such that $\f{1}{2}\|\rho^\sA - \pi_k^\sA\|_1 \geq 1/2$, we have $\rF(\rho^\sA, \pi_k^\sA) < 3/4$. 
Thus, estimating the squared fidelity within additive error of less than $1/8 \ (= (1 - 3/4)/2)$ requires $\Omega(k^{1/3})$ queries to the unitaries $U_\rho^{\sA\sB}$ and $U_{\Phi_k}^{\sA\sB}$ that prepare these purified states, as well as to their inverses.


As a result, for $\delta$ such that $\delta + 1/72 < 1/8$, i.e., for $\delta \in (0, 1/9)$, there exists a pair of rank-$k$ states $\rho^\sA$ and $\sigma^\sA$ for which implementing a quantum channel $\cT^\sB$ satisfying Eq.~\eqref{inteq:152} requires $q = \Omega(k^{1/3})$ queries to $U_\rho^{\sA\sB}$, $U_\sigma^{\sA\sB}$, and their inverses.
Noting that $\rF(\rho^\sA, \sigma^\sA) \geq \rF\big(\cT^\sB(\ketbra{\rho}{\rho}^{\sA\sB}), \ket{\sigma}^{\sA\sB}\big)$, we complete the derivation.


\end{proof}





